


Date: 2026-11-14. Branch \texttt{final} from tag \texttt{redesign-v1.1}. Key prefix \texttt{r3.}. Committed before any T1, T2 or T3 run. Nothing in this round tunes an algorithm: D1, quick retries, the severity table, UA-SCT, the Nessa-style policy, fixed 3 retries and the standard LoRaWAN-style retry rule are used as frozen in \texttt{redesign-v1.1} (\texttt{rd\_frozen\_A.json}, \texttt{rd\_frozen\_C.json}, \texttt{rd\_policy.py}, \texttt{world\_sim.py}, \texttt{s6\_early.py}, \texttt{ar\_r3.py}). Hardware data (\texttt{data/hardware}) is not used.

\subsection*{0. Why this round, and what was looked at before this file}

D1 lost its effect when moved from the custom simulator to ns-3, so a protocol claim that rests on one simulator is weak. This round uses link traces measured in real networks (TRACE\_DRIVEN), proofs and exhaustive checks (FORMAL), and real crash data (MEASURED\_SOFTWARE).

Looked at before this file was written, and nothing else: (1) the structure of the data files (column names, number of daily files, number of records of the Brno file: one device, 1,527 messages); (2) the README of LoED. \textbf{No loss rate, no autocorrelation and no replay result had been computed.} Every rule below that depends on the data (inclusion, group limits, event model) is fixed from the prompt and from the frozen code, not from the traces.

Labels: TRACE\_DRIVEN, FORMAL, MEASURED\_SOFTWARE, SIMULATED, ANALYTICAL, NOT\_RUN. No hardware label is used.

\subsection*{1. T1: trace-driven evaluation (TRACE\_DRIVEN)}

\subsubsection*{1.1 Data}

{\scriptsize\raggedright
\begin{longtable}{p{0.237\linewidth}p{0.237\linewidth}p{0.237\linewidth}p{0.237\linewidth}}
\toprule
Dataset & Source & Version and licence & Use \\
\midrule
\endhead
LoED, "The LoRaWAN at the Edge Dataset" (Bhatia, Breza, Marfievici, McCann; DATA '20; arXiv 2010.14211) & Zenodo record 4121430, file \texttt{LoED\_LoRaWAN\_at\_edge\_dataset.zip}, downloaded 2026-11-14, SHA-256 \texttt{8960cc3ad25b8bd3c99a5d9f371b191bae25e2bd5ca8b9f0cedd5fa4452cf3f3}; 188 daily CSV files, 1.5 GB unpacked & CC BY 4.0 (Zenodo record). Nine gateways in central London, 2019 to 2020. Columns used: time, device\_address, gateway, crc\_status, spreading\_factor, rssi, snr, mtype, fcnt & Primary (all hypotheses) \\
"Measurements of LoRaWAN Technology in Urban Scenarios: A Data Descriptor" (Data 2021, 6(6), 62; Brno University of Technology with Tampere University authors) & Zenodo record 6228358, downloaded 2026-11-14; file \texttt{exportMessages\_2019 LoRaWAN\_CRA\_1527\_original.json} & CC BY 4.0 (Zenodo record); the code mirror is MIT-licensed & Descriptive only: it holds one device (1,527 messages), so no split by device is possible and no hypothesis uses it \\
\bottomrule
\end{longtable}}

The raw files stay in \texttt{data/raw/} (git-ignored). Both licences allow the use; both are cited in the paper.

\subsubsection*{1.2 Loss sequences (rules fixed here)}

\begin{itemize}
\item LoED rows with \texttt{crc\_status = 1}, \texttt{device\_address != -1} and a valid \texttt{fcnt} are kept. Copies of one frame (same device, gateway and \texttt{fcnt} within 2 s: antenna or channel duplicates) are merged.
\item \textbf{Link = device} (network level): a frame counts as received if at least one of the nine gateways decoded it with a valid CRC. Gaps in the 16-bit frame counter of the union sequence are frames sent and not received by any gateway.
\item Consecutive received counters f, g give a gap of d = (g - f) mod 65536. For 1 <= d <= 100 the d - 1 frames in between are lost. d = 0 is a duplicate (ignored). d > 100, or a counter that decreases from below 60,000 (a reset), starts a new segment: it is not counted as loss (the device was off or silent).
\item Frames of one segment form the 0/1 reception sequence of the link (1 = received). Time stamps of lost frames are interpolated linearly between their received neighbours.
\item \textbf{Inclusion:} at least 300 frames (received plus lost) in total over its segments. \textbf{Groups by loss rate p (lost / total):} L: p < 0.10; M: 0.10 <= p < 0.40; H: p >= 0.40.
\item \textbf{Split:} by device, hashed: \texttt{int(sha256("t1|" + device\_address).hexdigest(), 16) \% 2}; 0 = DEVELOPMENT, 1 = HELD-OUT. Only HELD-OUT devices enter the hypotheses. The DEVELOPMENT half is used to check that the code runs and the invariants hold; no parameter is chosen on it (nothing is tuned).
\item Reported for the HELD-OUT half and for the DEVELOPMENT half: the per-device loss rate distribution (quantiles and group sizes), the lag-1 to lag-10 autocorrelation of the loss indicator, the burst-length distribution, and the ratio P(loss at t+1 | loss at t) / p (1 for independent losses; the custom simulator assumed independent losses), with RSSI and SNR summaries.
\end{itemize}

\subsubsection*{1.3 Replay (rules fixed here)}

Event: a Critical alert at crash instant 0 with deadline 10 s (V7 and V12 convention); delivery means reception at the base before the deadline (the transmission ends no later than 10 s). The alert is ready after the window of 1.8 s plus 0.05 s compute (\texttt{replay.COMPUTE\_S}, \texttt{config.POST\_S}) plus the policy's hold: 5 s for the severity table (P1) and for fixed 3 retries (P0), 0 s for the others (as frozen in the SCT machinery).

Start points: for each device, 300 start indices drawn uniformly (seed derived from the device hash and 42) from the positions of its segments that have at least 40 following frames. \textbf{Attempt j of an alert sees the reception state of frame s + j} (consecutive frames of the same device). The sub-second spacing of attempts therefore has no effect on the loss state; this is the trace's resolution (frames are seconds to hours apart) and is a stated limit: it favours no policy by design but removes any benefit of spreading retries beyond a fade, and it keeps the frame-to-frame burstiness.

Downlink rule: the ACK of a received attempt is received with probability 1 - p\_dl, drawn independently, where p\_dl = the link's own loss rate p (symmetric links, as in the custom simulator); the ACK matters for stopping, for D1's loss estimate and for airtime, not for delivery. Spreading-factor changes do not change reception in the replay (the trace fixes the link), but change airtime and the duty budget: this is conservative for policies that raise the SF. Base SF for all policies: 9 (the S-X convention); time on air from \texttt{lora.airtime\_ms}; duty limit 1\% of an hour (36 s of airtime in a sliding hour, nothing else on the node).

Policies (single-link reductions of the frozen ones; no crowd, no flash frame, no sharing of duty with other vehicles; stated):

{\scriptsize\raggedright
\begin{longtable}{p{0.475\linewidth}p{0.475\linewidth}}
\toprule
Name & Rule \\
\midrule
\endhead
D1 & frozen \texttt{D1Table(9)} and \texttt{BetaLoss} with the frozen parameters of \texttt{rd\_frozen\_A.json}, load bin 0; the Beta posterior is warmed with the 20 frames before the start (a round fails if the uplink or its ACK is lost); budget share 0.5 \\
quick & spacing 0.3 s after the listen window, SF + 1 after 2 rounds, up to 12 tries \\
P1 (severity table) & hold 5 s, retry at once after the listen window, up to 12 tries \\
UA-SCT (P4) & spacing 5 s for Critical, SF + 1 after 2 rounds, up to 12 tries \\
fixed 3 retries (P0) & hold 5 s, retry at once, 4 tries in total \\
standard LoRaWAN-style confirmed uplink (khan) & send at once, retry after the listen window plus a delay U(1, 3) s, up to 8 retries, no SF change \\
Nessa-style & frozen \texttt{k = 1}, \texttt{A = 1.0 s}; ACK timeout from the airtime, doubling per stage, a redundant copy after the initial retransmission of each stage, SF + 1 after each pair of retransmissions, up to 12 tries \\
\bottomrule
\end{longtable}}

Tamang-style slotted access needs a burst-size oracle and a crowd, so it is not part of the single-link replay (stated).

\subsubsection*{1.4 Hypotheses (all on HELD-OUT devices, Critical alerts)}

Device-level value of a policy: the mean over its 300 start points of the 0/1 delivery before the deadline. Group value: the mean over devices of the group. Paired bootstrap over devices (2,000 resamples, seed 42, the same resampled devices for every policy).

\begin{itemize}
\item \textbf{H-T1a (three tests: L, M, H).} In every group, D1 is not worse than the best of the six baselines of that group (the baseline with the highest group mean on the same HELD-OUT devices; the harsher choice) by more than 2 points: the upper 95\% limit of (best baseline - D1) is at most 0.02. One test per group; each has a one-sided bootstrap p-value (share of resamples with a difference above 0.02).
\item \textbf{H-T1b.} On HELD-OUT devices of group H (loss at least 40\%), D1 beats quick retries by at least 2 points: the point difference is at least 0.02 and the lower 95\% limit is above 0; one-sided p = share of resamples with a difference at most 0.
\item \textbf{H-T1c.} The provisional alert of R3 on the traces stays within the R3 bounds (\texttt{ar\_r3.py} logic: a 14-byte flash at the early-exit decision time of the stored per-event decisions of \texttt{s6\_early} at the 1\% target, seed 42, then the full alert or a CANCEL sent three times; at most 4 tries per message; the stage that follows uses the next frames): (a) the median paired difference (single-stage alert arrival minus provisional arrival), over crash events that alarm in both stages, is at least 0.8 s in each of the groups L, M and H; (c) extra transmit airtime is under 40\% and the event delivery of the two-stage scheme is not lower than that of the single-stage scheme by more than 2 points (lower 95\% limit of the difference at least -0.02) in each group. Decision times and decisions are drawn with replacement from the stored events, independent of the link. p-value: the largest of the clause p-values (intersection test).
\item Descriptive only (no hypothesis): the full table of all seven policies by group, with intervals; Severe alerts (30 s) are not evaluated.
\end{itemize}

\textbf{Kill criteria (T1).} Fewer than 100 HELD-OUT devices pass inclusion: T1 is reported as underpowered (not supported), and no group-level claim is made for a group with fewer than 20 devices (the test for that group is NOT\_RUN). Any failure of the invariant checks (a counted delivery after the deadline, a transmission above the duty limit, a policy using a frame twice) stops the run until the code is fixed, before the HELD-OUT run.

\subsubsection*{1.5 Limits to state}

The traces are uplink receptions, not ACK exchanges; downlink loss is modelled from the uplink by the rule above; the devices are not vehicles (they are mostly static sensors) and their duty cycle, payloads and spreading factors differ from the alert node; frame counters count frames sent, so silent periods of an idle gateway can inflate union loss; the frames are seconds to hours apart, so the loss state of retries seconds apart is not observed (section 1.3).

\subsection*{2. T2: formal guarantees (FORMAL)}

\subsubsection*{2.1 Security model (Tamarin 1.12.0 with Maude 3.5.1, installed in WSL; versions and model files are recorded with the run)}

Two models of the alert protocol: \textbf{M0}, the earlier design with one shared key for event-level ACKs (a holder of the key can forge an ACK for any vehicle); \textbf{M1}, the final design: a per-vehicle key derived from a master key and the vehicle id, the ACK bound to the vehicle id, the alert counter and the freshness value, GPS-time freshness with a window and a seen-set, and a counter fallback after a reboot (the model follows \texttt{PACKET\_FORMAT.md}, \texttt{PROTOCOL\_MODEL.md} and PREREGISTRATION\_V11 sections 4, R1 and R4). Time is modelled in epochs (a natural-number counter at the base; the window is the current and the previous epoch).

Lemmas with the expected result written here, before the runs:

{\scriptsize\raggedright
\begin{longtable}{p{0.237\linewidth}p{0.237\linewidth}p{0.237\linewidth}p{0.237\linewidth}}
\toprule
Lemma & Statement & Model & Expected \\
\midrule
\endhead
L1 alert authenticity & If the base accepts an alert of vehicle V, V sent it (injective agreement on vehicle, counter and epoch) & M1 & proved \\
L2 ACK authenticity per vehicle & If vehicle V accepts an ACK, the base sent it for V's alert (counter, epoch); no ACK built from the key of another vehicle W is accepted by V & M1 & proved \\
L3 no cross-vehicle suppression & An adversary that holds the key of vehicle W (a compromised node) cannot make a vehicle V != W accept an ACK for V's alert & M1 & proved \\
L3' cell-wide ACK suppression & The same adversary holding the shared key can make an honest vehicle stop its retries without the base having received its alert & M0 & attack found \\
L4 replay outside the window & An alert whose epoch is older than the previous epoch is never accepted & M1 & proved \\
L4b replay inside the window & An alert is accepted at most once while the base's seen-set is intact & M1 & proved \\
L5 car reboot & After a car reboot (counter reset, GPS time continuing) no alert accepted before the reboot is accepted again, and a fresh alert after the reboot is accepted (executability) & M1 & proved \\
L6 base reboot & After a base reboot (seen-set lost) a replay inside the window is not accepted & M1 & \textbf{expected refuted} (known residual: the model should find a replay inside the window after a base reboot; V11 gives this event a probability near 2e-5 per window) \\
\bottomrule
\end{longtable}}

H-T2sec: L1 to L5 are proved on M1, L3' finds the attack on M0 and L6 is refuted as expected. \textbf{Kill criterion:} any of L1 to L5 is refuted, or L3' is not found (the model is then wrong or incomplete and is repaired and rerun, with the change listed as a deviation). Tamarin's symbolic model assumes perfect cryptography; nothing here says the implementation is secure. The earlier work cited and not claimed over: Wesemeyer et al. (EuroS\&P 2020) and Dong et al. (ISSRE 2024).

\subsubsection*{2.2 Timed guarantees}

Tool: UPPAAL is the tool named in the task. Its licence is for academic use after a registration step that this run cannot complete; it is not used. The retry logic of each policy is a deterministic automaton over time on air, listen windows, holds and the duty blackout, so the check is an \textbf{exhaustive explicit-state enumeration of all loss patterns of the leading n attempts (n up to 14, 2\textasciicircum{}n patterns, 1 ms time resolution) with the frozen policy code} (own code, \texttt{r3\_t2\_timed.py}); where a timed-automata model checker (TChecker) can be built in WSL, the model of one policy is also checked in it and the two are compared. The label is FORMAL (exhaustive finite-model check); the tool is named in every table and is not UPPAAL.

Model: one Critical alert, crash at 0, deadline 10 s, ready time as in 1.3, base SF 9, duty limit 1\% in a sliding hour, listen window PROC + ACK airtime + margin after every transmission. For k = 0 to 6, the worst-case time to delivery under at most k consecutive lost attempts is the largest, over all loss patterns with no run of more than k consecutive lost uplinks, of the end time of the earliest received uplink (patterns with no reception within the deadline are counted as a missed deadline and listed). The duty check asks whether any reachable transmission exceeds the allowance.

\textbf{H-T2.} For every k in 0..6, D1's worst-case delivery time for Critical alerts is not longer than that of quick retries, and no reachable D1 transmission breaches the duty budget. \textbf{Kill criterion:} the checker finds a counterexample; it is then reported with the pattern. The severity-table, UA-SCT, fixed-3, Nessa-style and LoRaWAN-style entries are computed for comparison (the cut order allows dropping the severity table).

\subsection*{3. T3: two-stage alert from real crash data (MEASURED\_SOFTWARE)}

Data: the de-duplicated VZCrash test split (24,022 events) with the stored per-event decisions of seed 42 (\texttt{results/tables/s6\_events\_seed42\_a1.csv} and \texttt{a5.csv}: triggered, alarm of the early-exit rule E, alarm of the full-window forest F, decision time); thresholds were chosen on validation by \texttt{s6\_early}; no model is refitted here, and **the test split was already scored once by \texttt{s6\_early} and used again by R3 (\texttt{ar\_r3})**. T3 reuses the stored decisions; this is stated, not a new test of the models. NHTSA pulses: the pooled records of \texttt{nhtsa.py}.

Measured: (1) time from trigger to the provisional decision (the early-exit decision time) and to the full decision (1,800 ms), on triggered crash events that alarm in both stages; (2) the retraction share (provisional alarms that the full window does not confirm, divided by provisional alarms) at the 1\% and 5\% false-alarm targets; (3) on NHTSA pulses, the agreement between the severity class computed from the window truncated at each early-exit length (50, 100, 200, 400, 800 ms after the trigger) and the class from the full window (the severity code of \texttt{severity.py}, unchanged); weighted by the distribution of the early-exit decision times of crash events when available. Bootstrap over events (2,000 resamples, seed 42), 95\% intervals.

\begin{itemize}
\item \textbf{H-T3a.} The median provisional decision is at least 0.8 s earlier than the full decision (median of 1,800 ms minus the early decision time is at least 800 ms; the lower 95\% limit is at least 800 ms) at the 1\% target.
\item \textbf{H-T3b.} The retraction share is under 20\% at the 1\% target (the upper 95\% limit is below 0.20).
\item \textbf{H-T3c.} The end-to-end earlier arrival (median paired arrival difference of at least 0.8 s) holds in at least two of the three link models: the custom simulator (\texttt{ar.r3}, existing), ns-3, and the trace-driven model (H-T1c clause (a)). \textbf{ns-3 counts as not holding unless it is run:} the single-paper P1 has not run in this session and the ns-3 application has no two-stage alert, so ns-3 is NOT\_RUN unless a two-stage application is built in time.
\end{itemize}

\subsection*{4. T4: post-crash survivability budget (ANALYTICAL, supporting)}

From datasheet currents of the ESP32, the SX1278 and the GPS module: record each source; any value from a datasheet that was not opened is marked ASSUMPTION. Compute how many alert attempts per policy a stated backup capacitor or cell can supply and which policies can still meet the Critical deadline on backup energy. No hypothesis unless every input is verified.

\subsection*{5. Multiple testing and decision rule}

One Holm family of seven tests: H-T1a-L, H-T1a-M, H-T1a-H, H-T1b, H-T1c, H-T3a, H-T3b (at most 10). H-T2sec, H-T2 and H-T3c have no p-value (FORMAL checks and a count rule). A test with a p-value is called supported if its criterion holds at the point level and the Holm-adjusted p is below 0.05.

A protocol claim reaches the main text only if it holds in FORMAL or TRACE\_DRIVEN evidence, or in at least two of the three link models. The decision table (T6) lists every hypothesis with result, label, interval, nearest prior work (\texttt{NOVELTY\_MATRIX\_2021PLUS.md}), allowed wording and placement.

\subsection*{6. Seeds, code and files}

Seed 42 for every draw; no repeated-seed runs are needed because the traces are fixed. Code: \texttt{r3\_t1.py}, \texttt{r3\_t2\_*.py}, \texttt{r3\_t3.py}, \texttt{r3\_t4.py} and the Tamarin files in \texttt{formal/}. Results under \texttt{r3.} in \texttt{results/results.json} and \texttt{results/tables/r3\_*}.

